\paragraph{同代码槽位分离额外采样的解释。}
增加成员同时增加了代码和尝试次数。因此，我们将九个不同成员与九个逐字节相同的 \bare{} 槽位比较，每个槽位在每道题上执行三次。后者保留随机结果产生的机会，同时去掉程序专长。它是等能力参照，并不假定所有生成程序都应表现得像 \bare{}。

\paragraph{留出重复将发现与验证分开。}
\label{sec:calibration-test}
若用同一批结果挑选一道题的最佳程序，再用这些结果验证它，就会重复利用偶然成功。我们改为三折重复交叉拟合：每折用两次发现重复冻结逐任务成员映射与一个最佳固定成员，再用第三次重复测量二者的准确率差。跨折平均得到 $G_{\rm real}$ 与 $G_{\rm clone}$；固定比较对象由发现重复选出，不一定是 \bare{}。随后在两臂共同完整的任务上比较：
\begin{equation}
D=G_{\rm real}-G_{\rm clone}.
\label{eq:clone-difference}
\end{equation}
这里相减的是配对任务层面的增益。它检验：根据过去结果作出的选择，迁移至另一次执行时，不同程序是否比相同程序更有优势。由于过去结果来自同一批评估题，这属于重复诊断，不是全新任务上的路由检验。

\paragraph{缺失与灵敏度决定如何解释结果。}
完整任务要求每个成员的每次重复均有判决，未知结果保持缺失。描述量使用各臂自身完整任务，$D$ 使用交集。判断规则要求任一臂及交集的不完整任务均不超过 $10\%$、至少有 $30$ 个配对任务，并通过包括 $D>1$\,\% 在内的正效应检查。已知真值模拟检验现有重复次数能否识别交叉优势。这些检查将尚未判清与不存在的证据分开。附录~\ref{app:wp1r} 说明数据完整性；附录~\ref{app:statistics} 给出完整规则、并列处理、不确定性及敏感性分析。采集协议在采集前固定，采集中有操作性修订；v3 分析在采集之后修订。
